\chapter{Para qué sirve una empresa}\label{ch:value-creation}

% Alcance: el valor, el beneficio económico como señal, ROIC vs. WACC como la prueba maestra.
% Vínculo cruzado: LTV:CAC es la expresión por cliente del ROIC > WACC (\cref{ch:customer-economics}).

Cada métrica financiera del libro --- LTV:CAC, ROAS, NOPAT, valor terminal --- es un indicador
indirecto de una pregunta de fondo: ¿esta actividad genera retornos por encima del costo del
capital que consume? El beneficio económico responde esa pregunta directamente; el resto del libro
rastrea cómo aflora en cada dominio funcional.

\section{Beneficio económico}\label{sec:economic-profit}
% complejidad: 3

¿Ganamos más que un retorno justo sobre el capital que desplegamos? La utilidad contable ignora
el cargo al capital propio; el beneficio económico no. Esta brecha es donde la mayoría de las
empresas destruye valor silenciosamente. Los acreedores cobran intereses, y ese cargo aparece en
el estado de resultados. Los accionistas también tienen un costo de oportunidad --- podrían
invertir en el mercado en su lugar. La utilidad contable omite ese cargo. El beneficio económico
lo resta: un número positivo significa que la empresa ganó por encima de su costo de capital
total; uno negativo significa que habría sido mejor devolver el dinero a los inversionistas,
incluso si la utilidad contable fue positiva.

Sea \(\mathrm{NOPAT}\) el beneficio operativo neto después de impuestos (ingreso operativo menos
el escudo fiscal sobre las operaciones, antes de financiamiento), \(\mathrm{WACC}\) el costo de
capital ponderado, e \(\mathrm{IC}\) el capital invertido (capital propio más deuda neta desplegada
en la empresa). Entonces:
\begin{equation}\label{eq:economic-profit}
  \mathrm{EP}
    \;=\; \mathrm{NOPAT} - \mathrm{WACC}\times\mathrm{IC}
    \;=\; (\mathrm{ROIC}-\mathrm{WACC})\times\mathrm{IC}.
\end{equation}
La forma de la derecha separa el \emph{spread} (ROIC menos WACC) de la \emph{escala} (IC): un
\emph{spread} alto sobre una base de capital pequeña genera menos beneficio económico que un
\emph{spread} moderado sobre una grande. Ambas palancas importan.

La fórmula es válida solo cuando el NOPAT es limpio: únicamente partidas operativas, después de
impuestos, sin cargos extraordinarios. El WACC debe corresponder al riesgo de los flujos de caja
que se evalúan --- usar el WACC corporativo para una división nueva y especulativa subestima el
umbral requerido. El capital invertido debe capturar todo el capital consumido, incluidos los
arrendamientos operativos y los costos capitalizados de I+D cuando sean significativos. El
\emph{goodwill} de adquisiciones es especialmente peligroso: incluirlo infla IC y deprime el ROIC
tras el cierre, haciendo que una operación positiva parezca destructiva de valor.

\begin{example}
La empresa SaaS de referencia gana \(\mathrm{NOPAT}\approx\money{600000}\) sobre capital
invertido de \(\money{3000000}\), lo que da \(\mathrm{ROIC}=20\%\). Con \(\mathrm{WACC}=11\%\):
\[
  \mathrm{EP}
    = (\underbrace{0.20}_{\text{return earned}} - \underbrace{0.11}_{\text{return required}})
      \times \money{3000000}
    = \money{270000}.
\]
La empresa crea \money{270000} de valor por año --- el spread del \qty{9}{\percent}
sobre la base de capital. Duplicar IC manteniendo ROIC constante duplica el beneficio económico;
mantener ROIC mientras el WACC sube (p.\,ej.\ porque la beta aumenta) lo erosiona sin ningún
cambio operativo.
\end{example}

\begin{admonition}{Advertencia}
El \emph{cash flow} libre y el beneficio económico no son la misma cosa. El FCF puede ser positivo
mientras el beneficio económico es negativo: una empresa madura que monetiza una franquicia en
declive genera caja pero puede ganar por debajo de su costo de capital. A la inversa, una empresa
de alto crecimiento con FCF fuertemente negativo puede tener beneficio económico fuertemente
positivo si despliega capital a retornos muy superiores al WACC. Ambas métricas responden
preguntas distintas; ninguna sustituye a la otra.
\end{admonition}

El EVA (Economic Value Added) de Stern Stewart es una implementación registrada de
\cref{eq:economic-profit} con ajustes a las cifras contables reportadas (capitalizar I+D, revertir
reservas LIFO, etc.) para producir un número de capital económico. El MVA (Market Value Added)
acumula el valor presente de los beneficios económicos futuros esperados --- equivale a la prima
del valor de mercado sobre el capital invertido. La equivalencia DCF--beneficio económico es
exacta: el NPV de todos los beneficios económicos futuros, descontados al WACC, es igual al NPV
de todos los cash flows libres futuros \parencite{brealey2020}. Se utiliza el que sea más
fácil de estimar; coinciden cuando los supuestos son consistentes.

El spread ROIC--WACC de \cref{eq:economic-profit} reaparece a lo largo del libro: a nivel
de cliente como LTV:CAC (\cref{ch:customer-economics}), a nivel de empresa como el spread
de ventaja (\cref{ch:advantage-and-disruption}), y a nivel de portafolio en el DCF ponderado por
escenarios (\cref{ch:valuation}). \textcite{ross2022} y \textcite{brealey2020} desarrollan los
fundamentos de las finanzas corporativas; \textcite{berk6thcorporate} hace el puente hacia la
valoración basada en el mercado.

\section{ROIC vs.\ WACC --- la prueba maestra}\label{sec:roic-wacc}
% complejidad: 3

¿Está la empresa acumulando valor o consumiéndolo? Una razón --- ROIC en relación con WACC ---
responde esa pregunta para toda la empresa y para cualquier inversión propuesta dentro de ella.
El capital no es gratuito. Cuando una empresa gana más sobre cada dólar desplegado de lo que los
inversionistas exigen para suministrarlo, multiplica riqueza. Cuando gana menos, destruye riqueza
aunque sea rentable en sentido contable. La distancia entre ROIC y WACC, multiplicada por la base
de capital, es el beneficio económico (\cref{eq:economic-profit}).

El retorno sobre el capital invertido es simplemente la razón entre el retorno operativo y el
capital desplegado:
\begin{equation}\label{eq:roic}
  \mathrm{ROIC} \;=\; \frac{\mathrm{NOPAT}}{\mathrm{IC}}.
\end{equation}
El valor se crea cuando \(\mathrm{ROIC}>\mathrm{WACC}\), se destruye cuando
\(\mathrm{ROIC}<\mathrm{WACC}\), y es neutral cuando son iguales. El crecimiento amplifica
cualquiera de las dos condiciones: invertir más capital con \(\mathrm{ROIC}>\mathrm{WACC}\)
crea más beneficio económico de forma proporcional; invertir con \(\mathrm{ROIC}<\mathrm{WACC}\)
lo destruye más rápido.

El ROIC es sensible al denominador. Una empresa con activos ligeros (alta rotación de activos)
puede mostrar un ROIC alto con márgenes modestos; una con activos pesados debe ganar márgenes
gruesos solo para superar el WACC. Ambas lecturas son legítimas --- lo que importa es el
spread, no el nivel absoluto de ninguna de las dos métricas. El ROIC también fluctúa con
el ciclo económico; la cifra de un solo año puede ser engañosa; conviene usar un promedio
normalizado o del ciclo completo.

\begin{example}
Continuando con el ejemplo de referencia: \(\mathrm{ROIC}=\money{600000}/\money{3000000}=20\%\)
frente a \(\mathrm{WACC}=11\%\). El spread es del \qty{9}{\percent} --- fuertemente
creador de valor.

A nivel de cliente, \(\mathrm{LTV}:\mathrm{CAC}=\money{3150}/\money{600}=5.25\times\)
es la forma por cliente de la misma prueba: el retorno de por vida sobre el capital de adquisición
de clientes, comparado implícitamente con el umbral \(\mathrm{LTV}:\mathrm{CAC}=1\). Los detalles
de ese cálculo se encuentran en \cref{ch:customer-economics}; la equivalencia es que ambos miden
ROIC sobre un despliegue de capital definido a lo largo de su vida útil.
\end{example}

\begin{admonition}{Advertencia}
Optimizar el ROIC reduciendo el denominador es una vía confiable para destruir la empresa.
Recortar el capital invertido --- posponer el mantenimiento, privar al \emph{marketing} de
recursos, agotar el capital de trabajo --- eleva el ROIC medido en el corto plazo mientras
agota la base productiva que genera el NOPAT futuro. La métrica parece mejorar precisamente
a medida que el negocio se deteriora. Un ROIC al alza con NOPAT plano o decreciente es una señal
de alerta, no de éxito.
\end{admonition}

El EVA y el MVA (presentados en \cref{sec:economic-profit}) cuantifican el efecto acumulado de
los spreads sostenidos de ROIC--WACC. En un esquema de DCF, una empresa que sostiene
\(\mathrm{ROIC}>\mathrm{WACC}\) indefinidamente tiene un valor terminal por encima del valor en
libros; la que gana exactamente el WACC tiene un valor terminal igual al valor en libros. El
marco de valoración de McKinsey convierte esto en el principio organizador central
\parencite{damodaran2012}. En el nivel de adquisiciones, una operación crea valor solo si el ROIC
post-cierre sobre el capital total invertido (incluida la prima pagada) supera el WACC --- la
razón del rigoroso modelado de sinergias de integración antes de firmar.

El ROIC ancla el tratamiento de finanzas corporativas de \cref{ch:corporate-finance-core} y la
construcción completa del DCF en \cref{ch:valuation}. El costo de capital WACC se deriva desde
primeros principios en \cref{ch:risk-and-capital-structure}. \textcite{damodaran2012} ofrece el
tratamiento estándar de los impulsores de valor basados en ROIC.

\section{Valor para el accionista vs.\ valor para las partes interesadas}\label{sec:stakeholder-value}
% complejidad: 3

¿Para quién debe optimizar la empresa? La respuesta moldea la asignación de capital, los
incentivos ejecutivos y la estrategia --- y es genuinamente controvertida. La postura de primacía
del accionista sostiene que el deber fiduciario de la administración recae en los propietarios del
capital propio; cualquier beneficio para las partes interesadas es legítimo solo en la medida en
que sirva a los retornos de los accionistas. La postura de las partes interesadas sostiene que
clientes, empleados, proveedores y comunidades son co-creadores de valor cuyos intereses limitan
y, en última instancia, hacen posibles esos retornos. Ambas posiciones coinciden en más de lo que
el debate sugiere --- el valor sostenible para el accionista requiere clientes satisfechos y
empleados motivados --- pero divergen marcadamente en gobernanza, estructuras de remuneración y
concesiones aceptables.

La postura de primacía del accionista se formula con precisión en \textcite{ross2022}: maximizar
el valor presente de los cash flows libres futuros distribuibles al capital propio. La
réplica de las partes interesadas --- el marco de gestión estratégica de Freeman --- sostiene que
las relaciones con las partes interesadas son en sí mismas activos generadores de valor; ignorarlas
subestima los flujos de caja a largo plazo:
\begin{description}
  \item[Primacía del accionista] Los accionistas son reclamantes residuales; maximizar su riqueza
    es el objetivo inequívoco. Todas las demás relaciones con partes interesadas son insumos,
    valorados por la competencia de mercado. \parencite{brealey2020}
  \item[Teoría de las partes interesadas] La estrategia debe dar cuenta de las demandas de
    clientes, empleados, proveedores, comunidades y reguladores. Ignorar estas relaciones genera
    pasivos ocultos que se materializan como riesgo regulatorio, rotación de talento y defección
    de clientes.
  \item[Lógica dominante de servicio] \textcite{vargo2004service} reencuadra el valor como
    co-creado entre la empresa y sus clientes en el acto de servicio: el producto es un recurso
    operante, no un bien terminado. Bajo esta visión, el beneficio económico se acumula a partir
    de relaciones, no de transacciones.
\end{description}

La primacía del accionista falla empíricamente cuando las externalidades son grandes (p.\,ej.\
contaminación), cuando las relaciones son el activo principal (servicios profesionales,
plataformas), o cuando la presión sobre las utilidades de corto plazo destruye la capacidad de
inversión a largo plazo de la que depende el valor del capital propio. El modelo de partes
interesadas falla cuando la lista de partes constituyen es tan amplia que la administración no
tiene un objetivo de optimización claro y puede racionalizar cualquier decisión como si sirviera
a alguien.

\begin{example}
Aplicado a la empresa SaaS de referencia: el razonamiento de primacía del accionista apoya
aceptar cada dólar de adquisición de clientes cuyo LTV:CAC esperado supere el umbral del WACC
(\cref{eq:roic}). El razonamiento de partes interesadas añade: ¿el producto mejora realmente el
resultado del cliente y es la estructura de margen de contribución justa para los empleados que
lo entregan? En la práctica, el modelo SaaS alinea ambos: la renovación de la suscripción es
voluntaria, por lo que un LTV sostenido requiere entrega genuina de valor. El desalineamiento
suele aparecer en contratos de bloqueo empresarial o en precios opacos --- no en un SaaS de uso
transparente.
\end{example}

\begin{admonition}{Advertencia}
Tratar el reporte ESG como sustituto del análisis de beneficio económico confunde dos preguntas
distintas. Una puntuación ESG alta no implica \(\mathrm{ROIC}>\mathrm{WACC}\); una baja no
implica lo contrario. Cada una merece su propia métrica rigurosa: el beneficio económico para la
productividad del capital, la medición explícita del impacto social para las externalidades.
\end{admonition}

La doctrina Friedman --- «la responsabilidad social de los negocios es aumentar sus utilidades»
(\emph{New York Times}, 1970) --- sigue siendo el ancla estricta de la primacía del accionista.
La declaración de la Business Roundtable de 2019, firmada por 181 directores ejecutivos, rechazó
explícitamente la primacía del accionista en favor de compromisos con las partes interesadas. Si
esto representa un cambio genuino en la gobernanza o una señal reputacional está activamente
debatido; la literatura empírica encuentra evidencia mixta sobre si los compromisos con partes
interesadas mejoran o perjudican los retornos para los accionistas en horizontes largos
\parencite{kotler2021}.

La lógica dominante de servicio (\cref{sec:stakeholder-value}) y el marco de trabajos por hacer
(\cref{ch:segmentation-targeting-positioning}) convergen en una implicación común: la unidad de
valor no es el producto sino el resultado que logra el cliente. Las empresas que diseñan en torno
a resultados en vez de características tienden a ganar mayor \emph{retention} y menor CAC ---
lo que se refleja en un LTV:CAC más alto, que es el beneficio económico a nivel de cliente. Los
marcos son complementarios, no competidores, cuando el horizonte temporal es suficientemente largo.

La tensión entre el valor para el accionista y el valor para las partes interesadas reaparece en
gobernanza (\cref{ch:organizational-behavior}) y en la discusión de ética de precios de
\cref{ch:pricing}. \textcite{vargo2004service} provee la base de la lógica dominante de servicio;
\textcite{kotler2021} revisa las implicaciones para la estrategia de marketing.

\section{ESG como disciplina de asignación de capital}\label{sec:esg-capital}
% complejidad: 3

¿Cuándo pertenece un factor ambiental, social o de gobernanza a una decisión de asignación de
capital y no a un informe de sostenibilidad separado? La disciplina es rigurosa: un factor ESG
gana un lugar en el cálculo de asignación de capital únicamente si afecta de forma medible al
NOPAT, al WACC o al capital invertido --- las tres palancas de \cref{eq:roic} y
\cref{eq:economic-profit}. Los factores que no superan esa prueba pueden valer la pena divulgarlos,
pero pertenecen a un análisis distinto. Las normas IFRS S1 e S2 del ISSB formalizan esta lente de
materialidad financiera a nivel de reporte: S1 exige la divulgación de riesgos de sostenibilidad
que afecten el valor empresarial; S2 extiende eso específicamente al clima. El estándar europeo
de doble materialidad (ESRS bajo la CSRD) va más lejos y requiere divulgación cuando el impacto
en la sociedad es material \emph{incluso si} aún no se detecta ningún efecto financiero. Para
propósitos de asignación de capital, el subconjunto financieramente material es el operativo.

En la práctica, los factores ESG ingresan al modelo de capital a través de tres canales. Primero,
un riesgo reprecia el costo de capital: una empresa con riesgo de transición significativo (costos
regulatorios de carbono) o exposición física al clima encontrará que los prestamistas e
inversionistas de capital propio exigen una prima de riesgo mayor, elevando el WACC directamente.
Segundo, una ineficiencia operativa o un pasivo suprime el NOPAT: las fallas de salud y seguridad,
la interrupción de la cadena de suministro por eventos climáticos y las demandas por
discriminación reducen la utilidad operativa después de impuestos. Tercero, un factor infla el
capital invertido sin un retorno proporcional: los activos varados, las obligaciones de
remediación ambiental y los gastos de infraestructura diferidos aparecen todos en IC una vez que
la contabilidad se depura. Cuando alguno de estos canales no es trivial, el análisis ESG es
inseparable del análisis de beneficio económico.

\begin{example}
La empresa SaaS de referencia enfrenta un factor ESG que claramente supera el umbral de
materialidad financiera: una brecha salarial de género concentrada en el equipo de ingeniería. El
mecanismo es a través del NOPAT --- los ingenieros con salario inferior se van, lo que genera
búsquedas de reemplazo. Cada salida de un ingeniero de nivel medio tiene un costo estimado de
\money{80000} en reclutamiento, incorporación y pérdida de productividad. Un programa de auditoría
de equidad salarial y remediación cuesta \money{20000} y previene una salida al año.

La decisión es una aplicación estándar del NPV (\cref{eq:npv}). El beneficio anual es el costo de
reemplazo evitado de \money{80000}; el costo anual adicional es la auditoría de \money{20000};
el flujo de caja neto anual es \money{60000}. Tratando esto como una perpetuidad al WACC del
\qty{11}{\percent} de la empresa,
\[
  \mathrm{NPV} \;=\; \frac{\money{80000} - \money{20000}}{0.11} \;\approx\; \money{545000}.
\]
La inversión tiene un NPV fuertemente positivo --- no por óptica ESG, sino porque la prueba de
materialidad financiera se cumple. A la inversa, un programa de compensación de carbono que cuesta
\money{30000} al año y no tiene efecto medible en el NOPAT, el WACC o el IC no supera la prueba
y pertenece al presupuesto de RSE, no al plan de capital.
\end{example}

\begin{admonition}{Nota}
La advertencia de este capítulo de que una puntuación ESG alta no implica
\(\mathrm{ROIC}>\mathrm{WACC}\) --- y viceversa --- aplica directamente aquí. El ejemplo de
equidad salarial de arriba supera la prueba de materialidad financiera; la mayoría de las
puntuaciones ESG reportadas son agregados de factores, muchos de los cuales no la superan. La
disciplina consiste en desagregar: evaluar cada factor ESG frente a los tres canales del modelo de
capital (NOPAT, WACC, IC) antes de incluirlo en cualquier proyección.
\end{admonition}

Las relaciones con partes interesadas que \textcite{freeman2010strategic} identifica como activos
estratégicos son precisamente las que superan la prueba de materialidad financiera: empleados cuyo
compromiso sostiene el NOPAT, clientes cuya confianza sostiene la retention y limita el CAC,
y reguladores cuya buena voluntad evita el riesgo de licencia que inflaría el WACC. La lente de
co-creación de \textcite{vargo2004service} añade que la relación con el empleado es un recurso
operante --- no puede reemplazarse con una transacción; su deterioro reduce la capacidad productiva
que mide el NOPAT. Ambos marcos, leídos a través de la lente de asignación de capital, llegan a la
misma conclusión: ESG no es un ejercicio de reporte paralelo sino un insumo anterior al cálculo de
beneficio económico cuando se cumple la materialidad.

La infraestructura de datos ESG requerida para ejecutar este análisis a escala es un problema del
capítulo 17 (\cref{ch:data-and-ai-strategy}): qué métricas de sostenibilidad son de calidad para
la toma de decisiones, cómo se recolectan y cómo las evaluaciones de doble materialidad se
operacionalizan como un flujo de trabajo de datos. La supervisión del consejo sobre qué se divulga
y qué entra al plan de capital se aborda en \cref{ch:organizational-behavior}.
